\documentclass[11pt]{article}
\usepackage[margin=0.88in]{geometry}
\usepackage[T1]{fontenc}
\usepackage{amsmath,amssymb,amsthm}
\usepackage{graphicx}
\usepackage[hidelinks]{hyperref}
\usepackage{microtype}

\newtheorem{theorem}{Theorem}
\newtheorem{corollary}[theorem]{Corollary}
\newtheorem{remark}[theorem]{Remark}
\newcommand{\A}{\mathcal A}
\newcommand{\I}{\mathfrak k}
\newcommand{\ad}{\operatorname{ad}}

\title{No Adjoint-Induced Almost-Complex Structure\\
on a Proper Unital Algebra Quotient}
\author{fa\_banach\_001 / agent\_lane\_00}
\date{17 August 2026}

\begin{document}
\maketitle

\begin{center}
\fbox{\parbox{0.92\textwidth}{
\textbf{Claimed status: full negative solution of the sole open question in
Table 2 of arXiv:2410.22055, subject to human review.}
For the non-fixed crossed-product case, the table's question mark in the
adjoint column is a minus sign.  In fact the obstruction holds for every
proper ideal of every unital associative algebra.
}}
\end{center}

\section{Source question}

Let $\A$ be a unital $C^*$-algebra, let $\I$ be the proper closed ideal
used in the paper, and let $G/K$ be the associated homogeneous space.  Its
tangent space at the base point is naturally $\A/\I$.  An admissible
operator $J$ induces an almost-complex structure only if
\begin{equation}\label{eq:source-condition}
                 \operatorname{Ran}(J^2+I_{\A})\subseteq\I;
\end{equation}
this is Definition 2.5 of the source.  Equivalently, the operator induced
by $J$ on $\A/\I$ must square to minus the identity.

The source considers $J=\ad_d$, where $\ad_d(a)=[d,a]$.  Its final table
leaves open whether such an operator can induce an almost-complex structure
when
\[
 \A=C(X)\rtimes_\alpha\mathbb Z
\]
and the quotient corresponds to a non-fixed invariant set.

\begin{figure}[htbp]
\centering
\includegraphics[width=0.92\textwidth]{figures/open_question_crop.png}
\caption{The sole remaining almost-complex question in Table 2 on PDF
page 14 of arXiv:2410.22055.}
\end{figure}

\section{Answer}

\begin{theorem}[Unit obstruction]\label{thm:unit}
Let $\A$ be a unital associative algebra over $\mathbb R$ or $\mathbb C$,
let $\I\subsetneq\A$ be a proper two-sided ideal, and let $d\in\A$.
Then the inner derivation $J=\ad_d$ preserves $\I$, but it cannot satisfy
\[
                   \operatorname{Ran}(J^2+I_{\A})\subseteq\I.
\]
Equivalently, the induced inner derivation on the nonzero unital quotient
$\A/\I$ cannot be an almost-complex structure.
\end{theorem}

\subsection*{Idea of the proof}

An almost-complex structure has square $-I$, so it is automatically
invertible.  An inner derivation can never be invertible on a nonzero
unital quotient: it annihilates the unit.  In the source's range
formulation, this is visible by testing the single element $1$.

\begin{proof}
Because $\I$ is a two-sided ideal, $[d,\I]\subseteq\I$, so $J$ induces
\[
 \overline J=\ad_{d+\I}:\A/\I\longrightarrow\A/\I.
\]
The quotient is nonzero and unital because $\I$ is proper.  Its unit obeys
\[
                 \overline J(1+\I)=[d+\I,1+\I]=0.
\]
Thus $\overline J$ has a nonzero kernel and cannot satisfy
$\overline J^{,2}=-I_{\A/\I}$.  Directly in $\A$, the same contradiction
is the one-line calculation
\[
                  (J^2+I_{\A})(1)=1\notin\I.
\]
Hence the required range inclusion fails.
\end{proof}

\begin{corollary}[Resolution of the crossed-product question]
Let $\A=C(X)\rtimes_\alpha\mathbb Z$ and let $\widetilde{\mathfrak k}_Y$
be any proper ideal used in the source, including the ideal associated with
a non-fixed invariant set $Y$.  For no $d\in\A$ does
$\ad_d$ induce an almost-complex structure on the corresponding
homogeneous space $G/K$.
\end{corollary}

\begin{proof}
Apply Theorem~\ref{thm:unit} with
$\I=\widetilde{\mathfrak k}_Y$.
\end{proof}

\begin{remark}[Stronger derivation form]
The proof uses only that the operator annihilates the unit.  Therefore any
derivation $D:\A\to\A$ that preserves a proper ideal $\I$ fails to induce
an almost-complex structure on $\A/\I$: every derivation of a unital
real or complex algebra has $D(1)=0$.
\end{remark}

\section{Scope and checks}

The Nijenhuis condition is irrelevant to this obstruction.  Even if an
adjoint operator happens to have vanishing Nijenhuis torsion modulo the
ideal, it already fails the square condition required of an almost-complex
structure.  Consequently this result fills the question mark in Table 2;
it does \emph{not} classify the still-partially-open adjoint entries in
Table 1.

The source makes the same non-existence conclusion for the Calkin algebra
using Wintner's theorem on commutators.  The unit calculation above is
stronger for this purpose and works uniformly for the crossed-product case
left open in the source.

A bounded novelty check found only arXiv version 1, dated 29 October 2024.
The sole OpenAlex-indexed citing item was an unrelated chapter on the
infinite-dimensional Siegel disc; no later explicit resolution of the
table question was found.  The proof is algebraic and self-contained.

\begin{thebibliography}{9}
\bibitem{GLT}
T. Goli\'nski, G. Larotonda, and A. B. Tumpach,
\emph{Nijenhuis operators on homogeneous spaces related to
$C^*$-algebras}, arXiv:2410.22055v1 (2024); International Journal of
Geometric Methods in Modern Physics (2025), 2540040.
\end{thebibliography}

\end{document}
